\section{What the models say about interference}\label{sec:interp}

\paragraph{A distractor makes the target feel stronger, and more so the stronger it is}
Every model that predicts well gives the same average bimanual shift, about 4\um\ for the left hand and 6\um\ for the right, and the shift grows with the distractor amplitude.
This is what one expects if the vibrations on the two hands are partly summed before the judgement.
For the left hand the shift is proportional to the amplitude.
The right hand has, in addition, a shift that does not depend on amplitude; this could be a real difference between the hands or a shift that rises quickly at low amplitudes and then grows more slowly. With one monkey the two cannot be told apart. A shift of the form $b_2\,d^{\,p}$ would be the natural next variant to test.

\paragraph{Sensitivity does not change}
The slope of the psychometric curve is the same with and without the distractor.
The left-hand JND increase in Model~B comes from the bimanual lapses, which B can only absorb into a flatter curve.
The claim that the distractor reduces sensitivity should be dropped, or stated as not detected.

\paragraph{The extra errors look like answers driven by the distractor}
The bimanual ``high'' errors to weak targets are not spread evenly.
For the left hand they occur only when the distractor is strong (24--32\um, near or above the 28\um\ boundary), which suggests that on a few percent of trials a strong distractor decides the answer.
Their response times point the same way.
In the 24--32\um\ sessions these errors are only 35--50\,ms slower than correct answers (left hand 615 against 564\,ms, right 598 against 562\,ms).
Errors that cannot come from the distractor are much slower: ``low'' answers to strong targets take about 805\,ms against 546\,ms, and the right hand's weak-target errors in the 6--12\um\ sessions about 1,010 against 550\,ms.
A literal ``judges the distractor'' model does not fit, though. Sirius calls a 24--32\um\ vibration ``high'' about half of the time or less (its unimanual PSE is near 32\um), so judging such a distractor would add ``low'' answers to strong targets about as often as ``high'' answers to weak ones; strong targets drew few ``low'' answers in those sessions (Fig.~\ref{fig:ppc}, bottom). The errors are better described as occasional ``high'' answers pushed by a strong distractor than as reports of the distractor itself, and neither capture model fits better than D's amplitude-dependent lapse rates. The choices alone do not settle the mechanism; the response times are the more direct test (Section~\ref{sec:further}).

\paragraph{The criterion moved between sessions}
Over the four months Sirius became more reluctant to call a unimanual target ``high'': its PSE rose by about \critPseLeft\um\ for the left hand and \critPseRight\um\ for the right (amplitude version; slightly more in the time version).
Two readings fit the data equally well.
Sirius may have adjusted its criterion to the distractor: in sessions with strong distractors many bimanual targets feel stronger, and a monkey that keeps its overall rate of ``high'' answers steady would raise its criterion, as the left hand's numbers suggest.
Or the criterion simply drifted, for reasons unrelated to the distractor.
The design cannot separate the two, because the distractor amplitude rose with time.
Either way, the trend matters for the distractor effects: without it, part of the trend is mistaken for a presence effect.

\paragraph{Model choice}
Model~B spends its parameters on things the data do not need, such as separate slopes and biases for unimanual and bimanual trials in every session, and misses what they do need: lapse rates tied to the distractor and a trend in the criterion.
Model~E keeps what is needed (a bias, slope and lapse rates for each session, bimanual lapses that depend on the distractor, a presence offset and the criterion trend) and fixes the rest.
